% Fields Medal Laureates - Complete List (1936–2026)
% All 68 laureates organized by congress year
% Compile with: xelatex fields_medal_laureates.tex

\documentclass[a4paper,10pt]{article}

% ---- Fonts & Language ----
\usepackage{fontspec}
\usepackage{xeCJK}
\setCJKmainfont{PingFang SC}[BoldFont=PingFang SC Semibold]
\setCJKsansfont{PingFang SC}
\setmainfont{Times New Roman}[BoldFont=Times New Roman Bold]
\setsansfont{Helvetica Neue}

% ---- Layout ----
\usepackage[margin=1.2cm,top=1.8cm,bottom=1.8cm,landscape]{geometry}
\usepackage{longtable}
\usepackage{booktabs}
\usepackage{array}
\usepackage{xcolor}
\usepackage{colortbl}
\usepackage{fancyhdr}
\usepackage{tikz}
\usepackage{hyperref}
\usepackage{amssymb}

% ---- Colors ----
\definecolor{headerblue}{RGB}{41,65,122}
\definecolor{headerbg}{RGB}{220,230,245}
\definecolor{rowalt}{RGB}{245,248,252}
\definecolor{yearcolor}{RGB}{180,50,50}
\definecolor{fieldcolor}{RGB}{30,100,60}
\definecolor{mutedgray}{RGB}{100,100,100}
\definecolor{congressbg}{RGB}{255,245,225}
\definecolor{wolfcolor}{RGB}{180,120,0}

% ---- Page Style ----
\pagestyle{fancy}
\fancyhf{}
\fancyhead[L]{\small\sffamily\color{headerblue} 菲尔兹奖获奖者完整名录 · Fields Medal Laureates}
\fancyhead[R]{\small\sffamily\color{headerblue} 1936–2026 · 共 68 人}
\fancyfoot[C]{\small\sffamily\thepage}
\renewcommand{\headrulewidth}{0.4pt}

% ---- Custom Commands ----
\newcommand{\congressheader}[2]{%
  % #1 = year, #2 = host city
  \rowcolor{congressbg}
  \multicolumn{6}{l}{%
    \raisebox{0pt}[1.2em][0.4em]{%
      \bfseries\sffamily\color{yearcolor} \large #1 \normalsize
      \color{mutedgray}\enspace #2%
    }%
  } \\
  \midrule
}

\newcommand{\laureaterow}[6]{%
  % #1=Name, #2=Award(age), #3=Birth–Death, #4=Nationality, #5=Institution, #6=Field
  {\bfseries #1} & {\color{yearcolor}#2} & {#3} & {#4} & {\small #5} & {\small\color{fieldcolor}#6} \\[2pt]
}

% Wolf Prize badge: append after name for dual laureates
\newcommand{\wolfbadge}{\textsuperscript{\normalfont\scriptsize\sffamily\color{wolfcolor}$\bigstar$W}}

% ---- Document ----
\begin{document}

\begin{center}
  \vspace*{0.3em}
  {\Huge\bfseries\sffamily\color{headerblue} 菲尔兹奖获奖者名录}\\[6pt]
  {\Large\sffamily\color{headerblue} Fields Medal Laureates 1936–2026}\\[6pt]
  {\normalsize\color{mutedgray} 共 21 届 · 68 位获奖者 · 按国际数学家大会届次排列}\\[3pt]
  {\small\color{mutedgray} Data sourced from Wikipedia \& IMU}
\end{center}

\vspace{0.8em}

\begin{longtable}{
  >{\raggedright}m{3.8cm}
  >{\centering}m{2.0cm}
  >{\centering}m{2.2cm}
  >{\raggedright}m{2.8cm}
  >{\raggedright}m{5.5cm}
  >{\raggedright\arraybackslash}m{6.5cm}
}

\toprule
\rowcolor{headerbg}
\textbf{\sffamily 姓名 Name} &
\textbf{\sffamily 获奖（年龄）} &
\textbf{\sffamily 生卒} &
\textbf{\sffamily 国籍} &
\textbf{\sffamily 机构 Institution} &
\textbf{\sffamily 研究领域 Field} \\
\midrule
\endfirsthead

\toprule
\rowcolor{headerbg}
\textbf{\sffamily 姓名 Name} &
\textbf{\sffamily 获奖（年龄）} &
\textbf{\sffamily 生卒} &
\textbf{\sffamily 国籍} &
\textbf{\sffamily 机构 Institution} &
\textbf{\sffamily 研究领域 Field} \\
\midrule
\endhead

\midrule
\multicolumn{6}{r}{\small\itshape 续下页 \ldots} \\
\endfoot

\bottomrule
\endlastfoot

%% ============================================================
\congressheader{1936}{奥斯陆 Oslo}
%% ============================================================

\laureaterow{Lars Ahlfors\wolfbadge}{1936（29岁）}{1907–1996}{芬兰}{Harvard University}{复分析、黎曼曲面、拟共形映射}

\laureaterow{Jesse Douglas}{1936（39岁）}{1897–1965}{美国}{MIT / City College of New York}{变分法、极小曲面（Plateau 问题）}

\midrule

%% ============================================================
\congressheader{1950}{剑桥 Cambridge (USA)}
%% ============================================================

\laureaterow{Laurent Schwartz}{1950（35岁）}{1915–2002}{法国}{Université de Nancy}{分布理论、广义函数}

\laureaterow{Atle Selberg\wolfbadge}{1950（33岁）}{1917–2007}{挪威}{Institute for Advanced Study}{解析数论、筛法、素数定理}

\midrule

%% ============================================================
\congressheader{1954}{阿姆斯特丹 Amsterdam}
%% ============================================================

\laureaterow{Kunihiko Kodaira\wolfbadge}{1954（39岁）}{1915–1997}{日本}{Princeton University / 东京大学}{复流形、代数曲面分类、Hodge 理论}

\laureaterow{Jean-Pierre Serre\wolfbadge}{1954（27岁）}{1926–}{法国}{Collège de France}{代数拓扑、代数几何、数论}

\midrule

%% ============================================================
\congressheader{1958}{爱丁堡 Edinburgh}
%% ============================================================

\laureaterow{Klaus Roth}{1958（33岁）}{1925–2015}{英国}{University College London}{丢番图逼近、数论}

\laureaterow{René Thom}{1958（35岁）}{1923–2002}{法国}{Université de Strasbourg}{配边理论、微分拓扑、突变论}

\midrule

%% ============================================================
\congressheader{1962}{斯德哥尔摩 Stockholm}
%% ============================================================

\laureaterow{Lars Hörmander\wolfbadge}{1962（31岁）}{1931–2012}{瑞典}{Stockholm University / Lund University}{线性偏微分算子、微局部分析}

\laureaterow{John Milnor\wolfbadge}{1962（31岁）}{1931–}{美国}{Princeton University}{微分拓扑、异国球面、K 理论}

\midrule

%% ============================================================
\congressheader{1966}{莫斯科 Moscow}
%% ============================================================

\laureaterow{Michael Atiyah}{1966（37岁）}{1929–2019}{英国}{University of Oxford}{K 理论、Atiyah-Singer 指标定理}

\laureaterow{Paul Cohen}{1966（32岁）}{1934–2007}{美国}{Stanford University}{集合论、forcing 方法、连续统假设}

\laureaterow{Alexander Grothendieck}{1966（38岁）}{1928–2014}{法国（无国籍）}{IHÉS}{代数几何、概形、上同调}

\laureaterow{Stephen Smale\wolfbadge}{1966（36岁）}{1930–}{美国}{UC Berkeley}{微分拓扑、高维庞加莱猜想、动力系统}

\midrule

%% ============================================================
\congressheader{1970}{尼斯 Nice}
%% ============================================================

\laureaterow{Alan Baker}{1970（31岁）}{1939–2018}{英国}{University of Cambridge}{超越数论、线性形式对数}

\laureaterow{Heisuke Hironaka}{1970（39岁）}{1931–}{日本}{Harvard University}{代数几何、奇点消解}

\laureaterow{Sergei Novikov\wolfbadge}{1970（32岁）}{1938–2024}{苏联}{Moscow State University}{代数拓扑、Pontryagin 类}

\laureaterow{John G. Thompson\wolfbadge}{1970（38岁）}{1932–}{美国}{University of Cambridge}{有限群论、奇阶定理}

\midrule

%% ============================================================
\congressheader{1974}{温哥华 Vancouver}
%% ============================================================

\laureaterow{Enrico Bombieri}{1974（33岁）}{1940–}{意大利}{Institute for Advanced Study}{素数分布、多复变、极小曲面}

\laureaterow{David Mumford\wolfbadge}{1974（37岁）}{1937–}{美国}{Harvard University}{代数几何、模空间、几何不变量理论}

\midrule

%% ============================================================
\congressheader{1978}{赫尔辛基 Helsinki}
%% ============================================================

\laureaterow{Pierre Deligne\wolfbadge}{1978（33岁）}{1944–}{比利时}{IHÉS / IAS}{代数几何、Weil 猜想、Hodge 理论}

\laureaterow{Charles Fefferman\wolfbadge}{1978（29岁）}{1949–}{美国}{Princeton University}{多复变、调和分析}

\laureaterow{Grigory Margulis\wolfbadge}{1978（32岁）}{1946–}{苏联}{Yale University}{Lie 群格、刚性理论、遍历理论}

\laureaterow{Daniel Quillen}{1978（38岁）}{1940–2011}{美国}{MIT / University of Oxford}{代数 K 理论、同伦理论}

\midrule

%% ============================================================
\congressheader{1982}{华沙 Warsaw}
%% ============================================================

\laureaterow{Alain Connes}{1982（35岁）}{1947–}{法国}{IHÉS / Collège de France}{算子代数、非交换几何}

\laureaterow{William Thurston}{1982（35岁）}{1946–2012}{美国}{Princeton University}{三维流形几何化、双曲几何}

\laureaterow{Shing-Tung Yau \textnormal{丘成桐}\wolfbadge}{1982（33岁）}{1949–}{美国/中国}{Harvard University / 清华大学}{微分几何、Calabi 猜想、几何分析}

\midrule

%% ============================================================
\congressheader{1986}{伯克利 Berkeley}
%% ============================================================

\laureaterow{Simon Donaldson\wolfbadge}{1986（29岁）}{1957–}{英国}{University of Oxford / Imperial College}{四维流形、规范理论}

\laureaterow{Gerd Faltings}{1986（32岁）}{1954–}{德国}{Max Planck Institute, Bonn}{算术代数几何、Mordell 猜想}

\laureaterow{Michael Freedman}{1986（35岁）}{1951–}{美国}{UC San Diego / Microsoft}{四维拓扑、庞加莱猜想（拓扑版）}

\midrule

%% ============================================================
\congressheader{1990}{京都 Kyoto}
%% ============================================================

\laureaterow{Vladimir Drinfeld\wolfbadge}{1990（36岁）}{1954–}{苏联/美国}{University of Chicago}{Langlands 纲领、量子群}

\laureaterow{Vaughan Jones}{1990（38岁）}{1952–2020}{新西兰}{UC Berkeley}{von Neumann 代数、Jones 多项式}

\laureaterow{Shigefumi Mori \textnormal{森重文}}{1990（39岁）}{1951–}{日本}{Kyoto University}{代数几何、极小模型纲领}

\laureaterow{Edward Witten}{1990（38岁）}{1951–}{美国}{Institute for Advanced Study}{数学物理、量子场论与几何}

\midrule

%% ============================================================
\congressheader{1994}{苏黎世 Zürich}
%% ============================================================

\laureaterow{Jean Bourgain}{1994（39岁）}{1954–2018}{比利时}{IHÉS / IAS}{Banach 空间、调和分析、组合数论}

\laureaterow{Pierre-Louis Lions}{1994（38岁）}{1956–}{法国}{Collège de France / Université Paris-Dauphine}{非线性 PDE、最优控制、黏性解}

\laureaterow{Jean-Christophe Yoccoz}{1994（37岁）}{1957–2016}{法国}{Collège de France}{动力系统、小除数问题}

\laureaterow{Efim Zelmanov}{1994（39岁）}{1955–}{俄罗斯/美国}{UC San Diego}{代数、限制 Burnside 问题}

\midrule

%% ============================================================
\congressheader{1998}{柏林 Berlin}
%% ============================================================

\laureaterow{Richard Borcherds}{1998（38岁）}{1959–}{英国}{UC Berkeley}{顶点代数、怪兽月光猜想}

\laureaterow{Timothy Gowers}{1998（34岁）}{1963–}{英国}{University of Cambridge}{泛函分析、组合数学}

\laureaterow{Maxim Kontsevich}{1998（34岁）}{1964–}{俄罗斯/法国}{IHÉS}{数学物理、Witten 猜想、形变量子化}

\laureaterow{Curtis T. McMullen}{1998（40岁）}{1958–}{美国}{Harvard University}{复动力系统、Teichmüller 理论}

\midrule

%% ============================================================
\congressheader{2002}{北京 Beijing}
%% ============================================================

\laureaterow{Laurent Lafforgue}{2002（36岁）}{1966–}{法国}{IHÉS}{Langlands 对应（函数域）}

\laureaterow{Vladimir Voevodsky}{2002（36岁）}{1966–2017}{俄罗斯}{Institute for Advanced Study}{动机上同调、A$^1$ 同伦理论}

\midrule

%% ============================================================
\congressheader{2006}{马德里 Madrid}
%% ============================================================

\laureaterow{Andrei Okounkov}{2006（37岁）}{1969–}{俄罗斯/美国}{Princeton University / Columbia University}{表示论、随机分拆、代数几何}

\laureaterow{Grigori Perelman}{2006（40岁）}{1966–}{俄罗斯}{无（拒绝领奖）}{Ricci 流、庞加莱猜想}

\laureaterow{Terence Tao \textnormal{陶哲轩}}{2006（31岁）}{1975–}{澳大利亚/美国}{UCLA}{调和分析、PDE、加性数论}

\laureaterow{Wendelin Werner}{2006（38岁）}{1968–}{法国}{Université Paris-Sud / ETH Zürich}{概率论、SLE、二维布朗运动}

\midrule

%% ============================================================
\congressheader{2010}{海得拉巴 Hyderabad}
%% ============================================================

\laureaterow{Elon Lindenstrauss}{2010（40岁）}{1970–}{以色列}{Hebrew University of Jerusalem}{遍历理论、测度刚性、数论}

\laureaterow{Ngô Bảo Châu \textnormal{吴宝珠}}{2010（38岁）}{1972–}{越南/法国}{Université Paris-Sud / University of Chicago}{代数几何、基本引理}

\laureaterow{Stanislav Smirnov}{2010（40岁）}{1970–}{俄罗斯}{Université de Genève}{统计物理、渗流共形不变性}

\laureaterow{Cédric Villani}{2010（37岁）}{1973–}{法国}{Institut Henri Poincaré}{最优传输、Boltzmann 方程}

\midrule

%% ============================================================
\congressheader{2014}{首尔 Seoul}
%% ============================================================

\laureaterow{Artur Avila}{2014（35岁）}{1979–}{巴西/法国}{IMPA / Université Paris Diderot}{动力系统、谱理论、重整化}

\laureaterow{Manjul Bhargava}{2014（40岁）}{1974–}{加拿大/美国}{Princeton University}{数论、数的几何、椭圆曲线}

\laureaterow{Martin Hairer}{2014（39岁）}{1975–}{奥地利/英国}{University of Warwick / Imperial College}{随机 PDE、正则结构理论}

\laureaterow{Maryam Mirzakhani}{2014（37岁）}{1977–2017}{伊朗}{Stanford University}{黎曼曲面模空间、双曲几何}

\midrule

%% ============================================================
\congressheader{2018}{里约热内卢 Rio de Janeiro}
%% ============================================================

\laureaterow{Caucher Birkar}{2018（40岁）}{1978–}{英国（库尔德裔）}{University of Cambridge}{代数几何、Fano 簇有界性}

\laureaterow{Alessio Figalli}{2018（34岁）}{1984–}{意大利}{ETH Zürich}{最优传输、自由边界问题}

\laureaterow{Peter Scholze}{2018（30岁）}{1987–}{德国}{Universität Bonn / Max Planck Institute}{perfectoid 空间、p-adic 几何}

\laureaterow{Akshay Venkatesh}{2018（36岁）}{1981–}{澳大利亚}{Institute for Advanced Study}{数论、表示论、动力系统}

\midrule

%% ============================================================
\congressheader{2022}{赫尔辛基 Helsinki}
%% ============================================================

\laureaterow{Hugo Duminil-Copin}{2022（36岁）}{1985–}{法国}{IHÉS / Université de Genève}{统计物理、相变、概率论}

\laureaterow{June Huh \textnormal{许埈珥}}{2022（39岁）}{1983–}{韩国/美国}{Princeton University}{组合数学、Hodge 理论、代数几何}

\laureaterow{James Maynard}{2022（35岁）}{1987–}{英国}{University of Oxford}{解析数论、素数间隔}

\laureaterow{Maryna Viazovska}{2022（37岁）}{1984–}{乌克兰}{EPFL}{Fourier 分析、球堆积问题}

\midrule

%% ============================================================
\congressheader{2026}{费城 Philadelphia}
%% ============================================================

\laureaterow{Yu Deng \textnormal{邓煜}}{2026（35岁）}{1990–}{中国/美国}{University of Chicago}{PDE、Boltzmann 方程推导}

\laureaterow{John Pardon}{2026（37岁）}{1989–}{美国}{Princeton University}{辛几何、Fukaya 范畴}

\laureaterow{Jacob Tsimerman}{2026（38岁）}{1988–}{加拿大}{University of Toronto}{数论、o-minimality、算术几何}

\laureaterow{Hong Wang \textnormal{王虹}}{2026（35岁）}{1990–}{中国/美国}{NYU Courant Institute}{几何测度论、Kakeya 问题}

\end{longtable}

\vspace{1.5em}

\begin{center}
\begin{tikzpicture}
  \node[fill=headerbg, rounded corners=5pt, inner sep=10pt, text width=0.85\textwidth, align=center] {
    {\small\color{mutedgray}
    \textbf{统计}：21 届大会 · 68 位获奖者 · 覆盖 1936–2026 年 \\[3pt]
    {\color{wolfcolor}$\bigstar$\kern-0.1em W} = 同时获得沃尔夫数学奖（Wolf Prize in Mathematics）\\[3pt]
    首位女性获奖者：Maryam Mirzakhani（2014）\\
    拒绝领奖：Grigori Perelman（2006）\\
    最年轻获奖者：Jean-Pierre Serre（27 岁，1954）\\
    数据来源：Wikipedia \& International Mathematical Union (IMU)}
  };
\end{tikzpicture}
\end{center}

\end{document}
